\section{Conclusion}
\label{sec:conclusion}

The fill rule of a market-making backtest is an assumption about an event that
cannot be observed, and it is usually left untested. Order-level data from an
on-chain limit order book makes it testable. On real orders whose fate is known,
the two standard event rules miss the observed fill rate by factors of two to
six in opposite directions, and an exponential-intensity clock is set by its
calibration rather than by the data. One mechanism explains both event rules:
almost all of the queue in front of an order leaves by cancellation. A rule that
follows the queue, with the exchange's in-block processing order, reproduces the
observed fill rate exactly; in-block cancellation priority rescues a few percent
of orders, apparently from the most toxic flow.

For backtests, the fill rule alone moves a strategy's PnL by up to \tbd{} and
its Sharpe ratio by up to \tbd, and the ranking of strategies agrees across
rules with Kendall $\tau$ = \tbd. \tbd
